\documentclass{article}

\usepackage{amsmath,amssymb}
\usepackage{xurl}
\usepackage[hidelinks]{hyperref}
\setlength{\emergencystretch}{2em}

\input{command}

\title{Nilpotent Basis Elements in the Simple-Tensor Definition of
Cirac--P\'erez-Garc\'ia--Schuch--Verstraete}
\author{}
\date{2026-07-16}

\begin{document}
\maketitle
\gapnote{local-correction}{resolved}

\section{Source definition}

At line~822 of Ref.~\cite{Cirac2016MPDO_arXiv},
Cirac--P\'erez-Garc\'ia--Schuch--Verstraete define simplicity as follows:
``A tensor generating MPDO's is simple if none of the elements in a BNT is
nilpotent.'' Line~819 of the same source defines an element
$\mathcal B_k$ of the basis of normal tensors to be nilpotent when
$\tr_R(\rho^{(N)}(M_k))=0$ after tracing out $R\leq D$ sites, for every
sufficiently large $N$.

\section{Formal reading}

Let $A=\mathcal B_k$ be a normal BNT representative with virtual dimension
$D$, and define its physical-trace transfer matrix by
\begin{align}
    \mathcal T_A
        &= \sum_i A^{ii}.
    \label{eq:cpsvnil_physical_trace_transfer}
\end{align}
Here $\tr_R(\rho^{(N)}(A))=0$ means that the reduced density operator is zero
after tracing out $R$ consecutive sites. For physical multi-indices
$\mathbf i=(i_1,\ldots,i_{N-R})$ and
$\mathbf j=(j_1,\ldots,j_{N-R})$, its matrix entries are
\begin{align}
    \left(\tr_R\!\left(\rho^{(N)}(A)\right)\right)_{\mathbf i,\mathbf j}
        &=
        \tr\!\left(A^{i_1j_1}\cdots A^{i_{N-R}j_{N-R}}\mathcal T_A^R\right).
    \label{eq:cpsvnil_reduced_operator_entry}
\end{align}

Choose a sufficiently large $N$ for which the length-$(N-R)$ products of the
normal representative span the full virtual matrix algebra. The trace pairing
on that algebra is nondegenerate. Consequently,
\begin{align}
    \tr_R\!\left(\rho^{(N)}(A)\right)=0
    \text{ for every sufficiently large }N
        &\Longleftrightarrow
        \mathcal T_A^R=0.
    \label{eq:cpsvnil_partial_trace_transfer_equivalence}
\end{align}
Since the nilpotency index of a nilpotent $D\times D$ matrix is at most $D$,
the line-819 condition of Ref.~\cite{Cirac2016MPDO_arXiv} is precisely
nilpotency of $\mathcal T_A$.

The formal predicate \leanid{MPOTensor.IsSimple} existentially quantifies over
the positive physical blocking required at source line~815
\cite{Cirac2016MPDO_arXiv} and chooses a BNT sector presentation of the blocked
doubled-index tensor. Every representative has a positive number of copies,
every copy has nonzero weight, the representatives are normal and eventually
linearly independent, and the blocked tensor has the same positive-length
matrix product vectors as their weighted direct sum. The predicate records
non-nilpotency for every representative. The nonzero copy weights implement
the nonzero-coefficient convention recorded in
\path{docs/paper-gaps/cpsv16_bnt_uniqueness_zero_coefficient.tex}, without
imposing the unit-modulus normalization from source line~246
\cite{Cirac2016MPDO_arXiv}.

Positive-length nontriviality is a theorem rather than a defining assumption.
A nonzero copy-weight power sum cannot vanish eventually. At a sufficiently
large length where one sector coefficient is nonzero, eventual linear
independence makes the blocked matrix product vector nonzero. Physical
unblocking then gives
\leanid{MPOTensor.IsSimple.exists_mpo_ne_zero}. Isolated vanishing lengths
remain possible. Definition~4.7 of Ref.~\cite{Cirac2016MPDO_arXiv} does not
require the generated MPO to be nonzero at every positive chain length, and
the formal predicate imposes no such requirement.

The predicate \leanid{MPOTensor.IsSimpleCanonicalForm} is a separate,
normalized, fixed-representative condition on an already-blocked tensor. It is
used for the Appendix~C.2 argument of
Ref.~\cite{Cirac2016MPDO_arXiv} and includes a normalized sector decomposition
with the source's line-246 unit-weight convention. No equivalence with
\leanid{MPOTensor.IsSimple} is asserted.

The toric-code boundary tensor discussed after Theorem~4.9 of
Ref.~\cite{Cirac2016MPDO_arXiv} confirms this reading. Its second basis element
has the one-dimensional transfer matrix
\begin{align}
    \mathcal T
        &= 0.
    \label{eq:cpsvnil_toric_code_zero_transfer}
\end{align}
The tensor is therefore not simple, as required by the source's treatment of
that example outside the simple case.

The source states the predicate for ``a BNT'' of the blocked tensor and then
asserts that BNTs are unique up to gauge, phase, and permutation
\cite{Cirac2016MPDO_arXiv}. The theorem
\leanid{MPSTensor.IsBNTSectorPresentation.equiv_of_sameMPV₂Pos} matches two
presentations, and
\leanid{MPOTensor.bnt_basis_not_isNilpotent_iff} transports nilpotency through
the resulting permutation, virtual-dimension identification, gauge
transformation, and phase. Thus the choice of presentation is immaterial. The
predicate does not require a normalized horizontal canonical-form
presentation.

\section{Classification}

The formal predicate is complete for source lines~217--246 and
Definition~4.7 of Ref.~\cite{Cirac2016MPDO_arXiv}, under the identification in
Eq.~\ref{eq:cpsvnil_partial_trace_transfer_equivalence} of the line-819
partial-trace criterion with matrix nilpotency of the physical-trace transfer.
The predicate requires a nonempty BNT sector presentation, and positive-length
nontriviality follows from that presentation. It does not require nonvanishing
at every positive length, and the source imposes no such condition.

\bibliographystyle{alpha}
\bibliography{references}

\end{document}
